\documentclass[11pt]{article}
\usepackage[a4paper,margin=1in]{geometry}
\usepackage[T1]{fontenc}
\usepackage{lmodern,microtype}
\usepackage{amsmath,amssymb,amsthm,mathtools}
\usepackage{booktabs}
\usepackage[numbers,sort&compress]{natbib}
\usepackage[colorlinks=true,allcolors=blue]{hyperref}
\newtheorem{theorem}{Theorem}[section]
\newtheorem{proposition}[theorem]{Proposition}
\newtheorem{corollary}[theorem]{Corollary}
\theoremstyle{remark}\newtheorem{remark}[theorem]{Remark}
\title{From a Cubic H\'enon Channel to CM Frobenius Data:\\
An Exact Arithmetic Bridge and Its Scope Boundary}
\author{Liang Wang$^{*1}$\\
\small $^{1}$School of Artificial Intelligence and Automation, Huazhong University of Science and Technology\\
\small Wuhan 430074, P.R. China\\
\small $^*$Corresponding author}
\date{Preprint, August 2026}
\hypersetup{pdfauthor={Liang Wang},pdftitle={From a Cubic Henon Channel to CM Frobenius Data}}
\begin{document}\maketitle
\begin{abstract}
Artificial prime damping makes a three-channel H\'enon--Kummer block
determinant trace class but adds an external clock.  We replace that damping
by the minimal geometric carrier of cubic symmetry, the elliptic curve
\(E:y^2=x^3+1\).  Its order-three automorphism gives complex multiplication
by \(\mathbb Z[\zeta_3]\), and its local polynomial is
\(1-a_pT+pT^2\).  We prove directly that \(a_p=0\) for every good prime
\(p\equiv2\pmod3\), because the cube map is bijective in that residue class.
An exact certificate counts all good primes below 2,000, verifies every Hasse
bound, and independently enumerates \(E(\mathbb F_{p^2})\) at four primes.
This supplies a genuine arithmetic connection and intrinsic square-root
normalization.  It does not supply a H\'enon periodic-orbit determinant: the
resulting object is the CM elliptic \(L\)-function, with different local and
global analytic data from the Riemann \(\xi\)-function.
\end{abstract}

\section{Introduction}
The homogeneous area-preserving H\'enon route naturally produces a cubic
phase.  A direct three-channel Kummer representation is gauge trivial; a
nonfunctorial permutation survives locally but its raw prime divisor
accumulates; and scalar damping repairs convergence only by inserting a new
clock.  These failures suggest a change in mathematical object rather than
another weight adjustment.

The curve
\begin{equation}\label{eq:E}
 E:\quad y^2=x^3+1
\end{equation}
is the smallest smooth projective arithmetic object carrying the same cubic
symmetry.  Over a field containing \(\zeta_3\), the map
\((x,y)\mapsto(\zeta_3x,y)\) is an order-three endomorphism.  The curve has
\(j=0\) and complex multiplication by \(\mathbb Z[\zeta_3]\); see
\cite{Milne2026,Silverman2009}.  This paper asks what exact arithmetic data
the cubic channel obtains from this completion and where the H\'enon claim
must stop.

We obtain a positive local theorem, a two-implementation numerical
certificate, and a negative provenance conclusion.  The positive theorem is
intrinsic arithmetic.  The negative conclusion prevents an elliptic
\(L\)-function from being promoted coordinatewise to a Riemann or
H\'enon determinant.

\section{Local Frobenius factors}
For a prime \(p>3\), let
\[
 N_p=\#E(\mathbb F_p),\qquad a_p=p+1-N_p.
\]
The local zeta numerator is
\begin{equation}\label{eq:local}
 L_p(T)=1-a_pT+pT^2.
\end{equation}
Its roots have absolute value \(\sqrt p\) by the elliptic-curve Riemann
hypothesis, equivalently \(|a_p|\le2\sqrt p\).  This is intrinsic
square-root normalization, not a fitted factor.

\begin{theorem}[Inert-prime trace vanishing]\label{thm:inert}
If \(p>3\) and \(p\equiv2\pmod3\), then \(a_p=0\) and
\(L_p(T)=1+pT^2\).
\end{theorem}
\begin{proof}
Let \(\chi\) be the quadratic character of \(\mathbb F_p\), extended by
\(\chi(0)=0\).  Point counting gives
\[
 N_p=p+1+\sum_{x\in\mathbb F_p}\chi(x^3+1),
 \quad
 a_p=-\sum_x\chi(x^3+1).
\]
Because \(\gcd(3,p-1)=1\), the map \(x\mapsto x^3\) is a bijection of
\(\mathbb F_p\).  Therefore
\[
 \sum_x\chi(x^3+1)=\sum_u\chi(u+1)=0,
\]
and the claim follows.
\end{proof}

For split primes \(p\equiv1\pmod3\), the trace need not vanish.  Direct
counts give
\[
a_7=-4,\qquad a_{13}=2,\qquad a_{19}=8.
\]
Thus the cubic channel distinguishes the inert and split prime classes
without a fitted classifier.

\section{Quadratic-extension independent check}
Let \(\alpha_p,\beta_p\) be the roots of \eqref{eq:local}.  Then
\(\alpha_p+\beta_p=a_p\), \(\alpha_p\beta_p=p\), and
\[
 \alpha_p^2+\beta_p^2=a_p^2-2p.
\]
Hence
\begin{equation}\label{eq:p2}
 \#E(\mathbb F_{p^2})=p^2+1-(a_p^2-2p).
\end{equation}

The release code verifies \eqref{eq:p2} independently.  For each of
\(p=5,7,11,13\), it chooses a quadratic nonsquare \(d\), represents
\(\mathbb F_{p^2}\) as pairs in \(\mathbb F_p[u]/(u^2-d)\), and tests square
status for every \(x\) in the extension.  This computation does not reuse
the prime-field point list.

\begin{table}[h]
\centering
\caption{Frozen local sentinels.}
\begin{tabular}{rrrr}
\toprule
\(p\)&\(p\bmod3\)&\(a_p\)&\(L_p(T)\)\\\midrule
5&2&0&\(1+5T^2\)\\
7&1&-4&\(1+4T+7T^2\)\\
11&2&0&\(1+11T^2\)\\
13&1&2&\(1-2T+13T^2\)\\
19&1&8&\(1-8T+19T^2\)\\
\bottomrule
\end{tabular}
\end{table}

\section{What the bridge proves}
The construction solves the provenance defect of C40 at one level.  The
factor \(p\) in \eqref{eq:local} and the size \(\sqrt p\) of Frobenius are
forced by geometry.  The local traces are reproducible, and repetitions are
controlled by the recurrence
\(a_{p,r}=a_pa_{p,r-1}-pa_{p,r-2}\).

It does not solve the H\'enon identification.  The curve \eqref{eq:E} is
selected because it carries cubic symmetry, but no theorem maps primitive
H\'enon unstable orbits to its Frobenius classes.  Its Euler product is the
elliptic \(L\)-function of \(E\), not the Riemann zeta function.  Its degree,
bad factors, gamma factor, and nontrivial zeros are correspondingly
different.  This is a new arithmetic connection, not a Hilbert--P\'olya
realization.

\section{Evaluator audit}
Route A credits reproducible Frobenius data at A1, but only weakly because
H\'enon chronology is missing.  The classical elliptic Euler product earns
an analytic determinant label at A2.  CM and square-root normalization give
partial A3.  Cohomology is only an A4 formal hint relative to the H\'enon
map.  The strict tuple is
\[
(A1_{\rm WEAK},A2_{\rm ANALYTIC\ DET},
A3_{\rm PARTIAL\ ANALYTIC},A4_{\rm FORMAL\ HINT}).
\]

A limited Route-B audit stops at B1.  No H\'enon-derived Hilbert space,
densely defined operator, domain, or spectral parameter map has been given.
An arithmetic Euler product cannot substitute for those data.

\section{Conclusion}
The cubic H\'enon channel has a natural arithmetic completion, and that
completion is nontrivial: the \(j=0\) CM curve supplies exact inert/split
Frobenius behavior and intrinsic square-root normalization.  It also changes
the target.  The next decisive question is finite and local: can any virtual
combination of this \(H^1\) factor and Tate factors equal the Riemann local
factor without inserting an exact negative copy?  A three-prime comparison
will answer that question before any global operator is attempted.

\bibliographystyle{plainnat}\bibliography{references}
\end{document}
